% Carga de datos iniciales
\subsection{Carga de datos iniciales}\label{sec:6_5_datos_semilla}

\subsubsection{Datos semilla (dependencias, chatbot y personal)}
Con los microservicios en funcionamiento, sus tablas ya existen y se pueden cargar los datos iniciales. Estos comandos se ejecutan en el \textbf{VPS Base de Datos}. Primero se confirma que las tablas se crearon:
\begin{lstlisting}
# cd /apps/aplicaciones/defensoria/database
# bash podman-compose-bd.sh psql-defensoria
defensoria_db=# \dt
defensoria_db=# \q
\end{lstlisting}
\captura{back-10-tablas}{Tablas creadas por los microservicios en defensoria\_db}

Después se cargan los datos semilla:
\begin{lstlisting}
# bash podman-compose-bd.sh seed
\end{lstlisting}
\begin{table}[H]
\centering
\caption{Datos semilla}
\footnotesize
\begin{tabular}{p{4.2cm}p{4cm}p{6.2cm}}
\toprule
\textbf{Archivo} & \textbf{Tabla} & \textbf{Contenido} \\
\midrule
\texttt{01-dependencias\_seed.sql} & \texttt{dependencias} & Catálogo de unidades académicas y dependencias del IPN \\
\texttt{02-chatbot\_seed.sql} & \texttt{preguntas\_chatbot} & Preguntas y respuestas del asistente del portal \\
\texttt{03-personal\_test.sql} & \texttt{personal\_administrativo} & Una cuenta por cada rol del personal \\
\bottomrule
\end{tabular}
\end{table}
El script ejecuta cada archivo en una sola transacción (si algo falla no se carga nada), no vuelve a cargar \texttt{dependencias} ni \texttt{preguntas\_chatbot} si ya tienen datos, y avisa qué microservicio falta levantar si una tabla todavía no existe.
\captura{back-11-seed}{Resultado de la carga de datos semilla}

Las cuentas de recepcionistas y de abogados de la Subdefensoría se cargan aparte:
\begin{lstlisting}
# cd /apps/aplicaciones/defensoria/database-scripts/seeds
# podman exec -i defensoria-db psql -U postgres -d defensoria_db < 04-recepcionistas.sql
# podman exec -i defensoria-db psql -U postgres -d defensoria_db < 05-abogados.sql
\end{lstlisting}
Se comprueba el número de registros cargados:
\begin{lstlisting}
# podman exec defensoria-db psql -U postgres -d defensoria_db -c \
  "SELECT 'dependencias', count(*) FROM dependencias
   UNION ALL SELECT 'preguntas_chatbot', count(*) FROM preguntas_chatbot
   UNION ALL SELECT 'personal_administrativo', count(*) FROM personal_administrativo;"
\end{lstlisting}
\captura{back-12-conteo}{Número de registros cargados en cada tabla}

\nota{Las cuentas de los scripts semilla son de prueba y tienen contraseñas conocidas. Antes de entregar la plataforma se desactivan:\\
\texttt{UPDATE personal\_administrativo SET activo = false WHERE correo\_institucional LIKE '\%.test@ipn.mx';}}
